% !TEX root = ../thesis.tex

%=========================================================
\chapter{Experimental Methods} 
\label{ch:methods}
%=========================================================

    This chapter describes the experimental methods used to obtain reproducible measurements. The experiment comprised not only the superconducting nanostructure, but the complete measurement chain around it: sample fabrication, cryogenic mounting, filtered electrical wiring, room-temperature instrumentation, digital acquisition, and numerical evaluation. Each part determined the quality of the measured \textit{I--V} curves and derived \textit{dI--dV} spectra, and constrained the interpretation of extracted model parameters.

    I first describe the fabrication of mechanically controllable break junction samples in Sec.~\ref{section:sample-preparation}. The process followed established group experience, but I adapted the design and fabrication protocol to improve sample stability, contacting, and reproducibility. The discussion focuses on the steps that determined the mechanical and electrical properties of the final devices. Practical fabrication notes are provided in Appendix~\ref{appendix:sample-fabrication}.

    I then follow the measurement chain from the sample environment to the room-temperature electronics in Sec.~\ref{sec:methods:setup}. The discussion covers mechanical actuation, cold wiring, filtering, microwave coupling, voltage and current reconstruction, and grounding. It emphasizes the choices that affected electronic temperature, low-frequency noise, microwave irradiation, and the reliability of the measured \textit{I--V} curves. Further details of the cryogenic environment are given in Appendix~\ref{appendix:cryostat}.

    Finally, Sec.~\ref{sec:methods:digital} connects the physical setup to the data products analyzed in later chapters. I describe the Python based instrument control, structured measurement files, evaluation pipeline, offset correction, and numerical interfaces for tunnel regime BCS fits and finite transmission MAR analysis. Together, the fabrication, measurement environment, and digital infrastructure define the experimental basis for the results chapters.
